\documentclass[11pt]{article}
\usepackage[margin=1in]{geometry}
\usepackage[T1]{fontenc}
\usepackage{amsmath,amssymb,amsthm,mathtools}
\usepackage{booktabs,array}
\usepackage[hidelinks]{hyperref}
\usepackage[numbers,sort&compress]{natbib}
\newcommand{\C}{\mathbb{C}}
\newcommand{\E}{\mathbb{E}}
\newcommand{\tr}{\operatorname{tr}}
\newcommand{\Tr}{\operatorname{Tr}}
\newcommand{\Sym}{\operatorname{Sym}}
\newcommand{\ket}[1]{\lvert #1\rangle}
\newcommand{\bra}[1]{\langle #1\rvert}
\newcommand{\Ptor}{P_{\mathrm{tor}}}
\newtheorem{theorem}{Theorem}
\newtheorem{lemma}[theorem]{Lemma}
\newtheorem{corollary}[theorem]{Corollary}
\newtheorem{proposition}[theorem]{Proposition}
\theoremstyle{definition}
\newtheorem{definition}[theorem]{Definition}
\theoremstyle{remark}
\newtheorem{remark}[theorem]{Remark}

\title{Universal Cloning Is Optimal Exactly for Designs}
\author{Michael Weldon\\[2pt]
\small Independent Researcher, Las Vegas, NV, USA\\
\small ORCID: \href{https://orcid.org/0009-0001-6119-0560}{0009-0001-6119-0560}}
\date{Version 1 --- 18 October 2026}

\begin{document}
\maketitle

\begin{abstract}
Werner's universal cloner turns $T$ copies of an arbitrary pure state in $\C^d$ into $T+k$ approximate copies with
optimal global fidelity $D_T/D_{T+k}$, where $D_m=\dim\Sym^m(\C^d)$. When the input is drawn from a restricted
ensemble, a better cloner may exist. We show that, for optimal \emph{average global} fidelity, it exists unless the
ensemble is a design of order $T+k$: for every probability measure on pure states and all $T,k\ge1$, the
optimal value equals $D_T/D_{T+k}$ if and only if the measure is a complex projective $(T+k)$-design. Sufficiency
follows from known arguments; the new content is necessity, which we prove from semidefinite-programming
complementary slackness together with a rigidity lemma for operators acting on symmetric tensors. As corollaries,
Werner's cloner is optimal for the uniform ensemble of $n$-qubit stabilizer states exactly when $T+k\le3$ (for every
$n$), for symmetric informationally complete ensembles exactly when $T+k\le2$, and for complete sets of mutually
unbiased bases exactly when $T+k\le3$ in $d=2$ and $T+k\le2$ in $d\ge3$. We give exact rational certificates for
the first strict gaps, and we contrast the result with structured state families, where matching moments is neither
necessary nor sufficient for optimality.
\end{abstract}

\section{Introduction}
Approximate cloning was introduced by Bu\v{z}ek and Hillery~\cite{BuzekHillery1996} and optimized for universal and
state-dependent qubit cloners by Bru\ss\ \emph{et al.}~\cite{Bruss1998}. Werner~\cite{Werner1998} solved the universal
case in general: from $T$ copies of an
unknown state in $\C^d$, the best global fidelity of $T+k$ output copies is $D_T/D_{T+k}$, attained by a covariant
channel that achieves this fidelity on every input. When the input is known to come from a restricted family, more
can be done, and determining the optimum for structured families is an active topic: phase states and other resolving
orbits~\cite{Herasymenko2026}, orthogonal orbits~\cite{CwiklinskiStudzinski2026}, finite sets in the asymptotic
regime~\cite{ChiribellaYang2014}, mixed states~\cite{Fanizza2026}, and stabilizer states~\cite{Bansal2026}, for which
Herasymenko \emph{et al.} note that the optimal cloning channel remains open~\cite{Herasymenko2026}.

A natural first question for any family is whether it is easier to clone than an arbitrary state at all. For
average global fidelity, the universal value is always achievable, so the question is when it is also optimal. It is
known that sufficiently uniform ensembles cannot be cloned better than the universal value: Molina, Vidick and
Watrous show this for $1\to2$ cloning of complex projective 3-designs~\cite{MVW2012}, and Kay notes the same
3-design condition for $1\to2$ cloning~\cite{Kay2014}. In measurement-based settings, designs give exact
characterizations: Scott characterizes the optimal measurement-based cloners~\cite[Thm.~20]{Scott2006}, and Kim
characterizes the ensembles attaining minimal quantumness~\cite{Kim2006}.

\paragraph{Result.} We prove that for coherent cloning the exact condition is a design of order $T+k$, in both
directions (Theorem~\ref{thm:main}). At $1\to2$ this is a 2-design, sharper than the 3-design sufficient condition
of~\cite{MVW2012}. The ``if'' direction is a short known argument; the ``only if'' direction is new as far as we can
determine (Section~\ref{sec:discussion}). Applied to standard ensembles, it settles exactly when Werner's cloner is
optimal for stabilizer, SIC and MUB ensembles (Corollary~\ref{cor:ensembles}), which addresses an average-global-fidelity
version of the stabilizer question raised in~\cite{Herasymenko2026} (who measure cloning precision by trace distance):
the universal cloner is optimal up to output order three and strictly suboptimal beyond it.

\section{Setting}\label{sec:setting}
Fix $d\ge2$, $T\ge1$, $k\ge1$, and write $n=T+k$, $D_m=\dim\Sym^m(\C^d)=\binom{d+m-1}{m}$, and $\Pi_m$ for the
projector onto $\Sym^m(\C^d)\subseteq(\C^d)^{\otimes m}$. An \emph{ensemble} is a probability measure $\mu$ on unit
vectors of $\C^d$ (phases are irrelevant); it may be finite with arbitrary weights, or continuous. Its $m$-th moment is
\[
M_m=\int(\ket\psi\!\bra\psi)^{\otimes m}\,d\mu(\psi).
\]
\begin{definition}
$\mu$ is a (complex projective, weighted) $t$-design if $M_t=\Pi_t/D_t$.
\end{definition}
A $t$-design is an $s$-design for all $s\le t$ (take partial traces). Equivalently~\cite[Thm.~6]{Scott2006}, the
frame potential satisfies $\iint|\langle\psi|\phi\rangle|^{2t}\,d\mu\,d\mu\ge1/D_t$, with equality if and only if $\mu$
is a $t$-design. Note that the design condition refers to the measure $\mu$ the cloner faces, not to its support.

A cloner is a channel (completely positive, trace preserving) $\Phi$ from operators on $\Sym^T(\C^d)$ to operators on
$\Sym^n(\C^d)$. Restricting the output to $\Sym^n$ loses nothing, since the targets lie there. Its average global
fidelity is
\[
F_\mu(\Phi)=\int\bra{\psi^{\otimes n}}\Phi\big(\ket{\psi^{\otimes T}}\!\bra{\psi^{\otimes T}}\big)\ket{\psi^{\otimes n}}\,d\mu(\psi),
\qquad P_\mu(T,k)=\max_\Phi F_\mu(\Phi).
\]
Werner's cloner
\[
\Phi_U(\rho)=c\,\Pi_n(\rho\otimes I^{\otimes k})\Pi_n,\qquad c=D_T/D_n,
\]
is a channel and satisfies $\bra{\psi^{\otimes n}}\Phi_U(\psi^{\otimes T})\ket{\psi^{\otimes n}}=c$ for every unit
$\psi$~\cite{Werner1998}. Hence
\begin{equation}\label{eq:lower}
P_\mu(T,k)\ge D_T/D_n\qquad\text{for every }\mu.
\end{equation}

\paragraph{Semidefinite program.} Fix an orthonormal basis $\{e_a\}_{a=1}^{D_T}$ of $\Sym^T(\C^d)$ and write
$s_a(\psi)=\langle e_a|\psi^{\otimes T}\rangle$. We use the Choi operator
$J=\sum_{a,b}\Phi(\ket{e_a}\!\bra{e_b})\otimes\ket a\!\bra b$ on $\Sym^n(\C^d)\otimes\C^{D_T}$, with the input register
last. Then $F_\mu(\Phi)=\tr(JX)$, where
\[
X=\int\ket v\!\bra v\,d\mu,\qquad \ket v=\ket{\psi^{\otimes n}}\otimes\textstyle\sum_a\overline{s_a(\psi)}\ket a .
\]
The cloning problem is the semidefinite program
\begin{equation}\label{eq:sdp}
P_\mu(T,k)=\max\{\tr(JX):J\succeq0,\ \Tr_{\rm out}J=I\}=\min\{\tr Y:I\otimes Y\succeq X\},
\end{equation}
where strong duality and attainment hold because $J=I/D_n$ is strictly feasible (for this formulation of
cloning as a semidefinite program see~\cite[Secs.~2--3]{MVW2012}). All transposes below refer to this
fixed basis and convention. Werner's channel has Choi operator
$J_U=c\sum_{a,b}\Pi_n(\ket{e_a}\!\bra{e_b}\otimes I)\Pi_n\otimes\ket a\!\bra b$.

\section{Main theorem}\label{sec:main}
\begin{theorem}\label{thm:main}
For every ensemble $\mu$ on pure states of $\C^d$ and all $T,k\ge1$,
\[
P_\mu(T,k)=\frac{D_T}{D_{T+k}}\quad\Longleftrightarrow\quad \mu\text{ is a }(T+k)\text{-design}.
\]
\end{theorem}

\begin{proof}[Proof of ``$\Leftarrow$'']
For any channel $\Phi$ and unit $\psi$, $\ket{\psi^{\otimes T}}\!\bra{\psi^{\otimes T}}\preceq\Pi_T$ and $\Phi$ is
positive, so $\Phi(\psi^{\otimes T})\preceq\Phi(\Pi_T)$. Integrating against $\psi^{\otimes n}$ gives
$F_\mu(\Phi)\le\tr[M_n\Phi(\Pi_T)]=\tr\Phi(\Pi_T)/D_n=D_T/D_n$ when $M_n=\Pi_n/D_n$, using that $\Phi$ maps into
$\Sym^n$ and preserves trace. Together with~\eqref{eq:lower} this gives equality. (For uniform finite ensembles this is
the moment bound of~\cite[Lemma~5]{Weldon2026}; for $T=k=1$ and 3-designs it is~\cite[Prop.~4]{MVW2012}.)
\end{proof}

The converse rests on two elementary lemmas about symmetric tensors.

\begin{lemma}\label{lem:A}
Let $k\ge1$, $w\in\Sym^T(\C^d)$ and $\phi\in\C^d$ a unit vector. If $w\otimes\phi^{\otimes k}\in\Sym^{T+k}(\C^d)$, then
$w\in\C\,\phi^{\otimes T}$.
\end{lemma}
\begin{proof}
Choose an orthonormal basis with $\phi=e_1$ and expand $w=\sum w_{i_1\cdots i_T}e_{i_1}\otimes\cdots\otimes e_{i_T}$.
The coefficient of $e_{i_1}\otimes\cdots\otimes e_{i_T}\otimes e_1^{\otimes k}$ in $w\otimes e_1^{\otimes k}$ is
$w_{i_1\cdots i_T}$. If some $i_m\ne1$, transposing slot $m$ with the last slot maps this basis tensor to one whose
last index is $i_m\ne1$; that coefficient is $0$ in $w\otimes e_1^{\otimes k}$. Full symmetry forces the two
coefficients to agree, so $w_{i_1\cdots i_T}=0$ unless all indices equal $1$.
\end{proof}

\begin{lemma}\label{lem:B}
Let $Z$ be a Hermitian operator on $\Sym^T(\C^d)$ such that every $\phi^{\otimes T}$ ($\phi$ a unit vector) is an
eigenvector of $Z$. Then $Z$ is a multiple of $\Pi_T$.
\end{lemma}
\begin{proof}
Let $Z\phi^{\otimes T}=\lambda(\phi)\phi^{\otimes T}$. Eigenvectors of a Hermitian operator with distinct eigenvalues
are orthogonal, and $\langle\phi^{\otimes T}|\chi^{\otimes T}\rangle=\langle\phi|\chi\rangle^T$. For unit vectors
$\phi,\chi$, pick a unit $\eta$ with $\langle\phi|\eta\rangle\ne0\ne\langle\chi|\eta\rangle$; then
$\lambda(\phi)=\lambda(\eta)=\lambda(\chi)$. So $\lambda$ is constant, and since the vectors $\phi^{\otimes T}$ span
$\Sym^T(\C^d)$, $Z=\lambda\Pi_T$.
\end{proof}

\begin{proof}[Proof of ``$\Rightarrow$'' in Theorem~\ref{thm:main}]
Assume $P_\mu(T,k)=c=D_T/D_n$. By~\eqref{eq:lower} Werner's channel is optimal, so $J_U$ is an optimal primal point.
Let $Y$ be an optimal dual point; $\tr Y=c$ and $Y=Y^\dagger$. Since $\tr[(I\otimes Y-X)J_U]=\tr Y-c=0$ and both
factors are positive semidefinite,
\begin{equation}\label{eq:cs}
(I\otimes Y-X)\,J_U=0 .
\end{equation}

\emph{Step 1: evaluating~\eqref{eq:cs}.} Define the operator $Z=\sum_{a,b}Y_{ba}\ket{e_a}\!\bra{e_b}$ on $\Sym^T(\C^d)$
(the transpose of $Y$ in the basis $\{e_a\}$); it is Hermitian and $\Pi_TZ=Z\Pi_T=Z$. Since
$\sum_aY_{a'a}\ket{e_a}=Z\ket{e_{a'}}$,
\[
(I\otimes Y)J_U=c\sum_{a,b}\Pi_n(Z\otimes I)(\ket{e_a}\!\bra{e_b}\otimes I)\Pi_n\otimes\ket a\!\bra b .
\]
For the second term, $\bra{\psi^{\otimes n}}\Pi_n=\bra{\psi^{\otimes n}}$ and $\sum_a s_a\overline{s_a}=1$ give
$\bra vJ_U=c\sum_b(\bra{e_b}\otimes\bra{\psi^{\otimes k}})\Pi_n\otimes\bra b$, and
$\int\overline{s_a}\,\ket{\psi^{\otimes n}}\bra{\psi^{\otimes k}}\,d\mu=M_n(\ket{e_a}\otimes I)$, so
\[
XJ_U=c\sum_{a,b}M_n(\ket{e_a}\!\bra{e_b}\otimes I)\Pi_n\otimes\ket a\!\bra b .
\]
Hence~\eqref{eq:cs} is equivalent to
$\big[\Pi_n(Z\otimes I)-M_n\big](\ket{e_a}\!\bra{e_b}\otimes I)\Pi_n=0$ for all $a,b$.

\emph{Step 2: a one-sided operator identity.} The ranges of $(\bra{e_b}\otimes I)\Pi_n$, $b=1,\dots,D_T$, lie in
$\Sym^k(\C^d)$ and span it: applied to $\psi^{\otimes n}$ they give $s_b(\psi)\,\psi^{\otimes k}$, and some
$s_b(\psi)\ne0$. With $Q=\Pi_T\otimes\Pi_k$, Step~1 is therefore equivalent to
\[
\Pi_n(Z\otimes\Pi_k)\,Q=M_n\,Q .
\]
Take adjoints. We have $Z^\dagger=Z$, $Q(Z\otimes\Pi_k)=Z\otimes\Pi_k$ (as $\Pi_TZ=Z$), $Q\Pi_n=\Pi_n$ and
$QM_n=M_n$ (as $\Sym^n\subseteq\Sym^T\otimes\Sym^k$ and $M_n$ is supported on $\Sym^n$). Therefore
\begin{equation}\label{eq:onesided}
(Z\otimes\Pi_k)\,\Pi_n=M_n .
\end{equation}

\emph{Step 3: rigidity.} Apply~\eqref{eq:onesided} to $\phi^{\otimes n}$ for an arbitrary unit vector $\phi$:
$(Z\phi^{\otimes T})\otimes\phi^{\otimes k}=M_n\phi^{\otimes n}\in\Sym^n(\C^d)$. By Lemma~\ref{lem:A},
$Z\phi^{\otimes T}\in\C\phi^{\otimes T}$ for every $\phi$, and by Lemma~\ref{lem:B}, $Z=\lambda\Pi_T$.

\emph{Step 4: normalization.} $\lambda D_T=\tr Z=\tr Y=c$, so $\lambda=1/D_n$, and~\eqref{eq:onesided} gives
$M_n=\lambda(\Pi_T\otimes\Pi_k)\Pi_n=\Pi_n/D_n$.
\end{proof}

\begin{remark}\label{rem:scope}
(i) The proof uses neither finiteness nor uniform weights, and it does not use Werner's uniqueness theorem; it uses
only that Werner's channel attains $c$ on every pure state. (ii) The hypothesis $k\ge1$ is necessary: for $k=0$ the
identity channel has fidelity $1$ for every ensemble. (iii) The statement concerns \emph{average global} fidelity.
For worst-case global fidelity the ``if'' direction persists (the worst-case optimum lies between $c$ and the average
optimum), but we make no claim about necessity. For single-copy (local) fidelity the picture differs: for a
continuous orthogonal orbit that forms a 2-design, the averaged $1\to2$ local-fidelity problem coincides with
universal cloning~\cite{CwiklinskiStudzinski2026}, and three equatorial qubit states reproduce phase-covariant
single-copy cloning~\cite{Jing2012}. (iv) The output order $T+k$ is sharp: the
objective $X$ depends on the moment $M_{2T+k}$, yet only $M_{T+k}$ decides equality with the universal value.
\end{remark}

\section{Consequences for standard ensembles}\label{sec:cor}
\begin{corollary}\label{cor:ensembles}
Werner's universal cloner is optimal for average global fidelity, i.e.\ $P_\mu(T,k)=D_T/D_{T+k}$, exactly in the
following cases.
\begin{enumerate}
\item[(a)] $\mu$ uniform on all $n$-qubit stabilizer states, any $n\ge1$: if and only if $T+k\le3$.
\item[(b)] $\mu$ uniform on a SIC (symmetric informationally complete set of $d^2$ states) in any $d\ge2$: if and only if
$T+k\le2$.
\item[(c)] $\mu$ uniform on a complete set of $d+1$ mutually unbiased bases: if and only if $T+k\le3$ when $d=2$, and
$T+k\le2$ when $d\ge3$.
\end{enumerate}
\end{corollary}
\begin{proof}
By Theorem~\ref{thm:main} it suffices to know the largest $t$ for which each ensemble is a $t$-design.
(a) Uniform stabilizer states form a 3-design and not a 4-design for every number of qubits~\cite{KuengGross2015}.
(b) SICs are 2-designs~\cite{Renes2004}. With overlaps $|\langle\psi|\phi\rangle|^2=1/(d+1)$ between distinct elements,
the order-3 frame potential is $\big(1+(d^2-1)(d+1)^{-3}\big)/d^2=(d+3)/(d(d+1)^2)$, which exceeds
$1/D_3=6/(d(d+1)(d+2))$ for every $d\ge2$ (the difference has the sign of $d^2-d$). So no SIC is a 3-design.
(c) Complete sets of MUBs are 2-designs~\cite{KlappeneckerRoetteler2005}. With overlaps $0$ within a basis and $1/d$
across bases, the order-$t$ frame potential is $\big(1+d^2\,d^{-t}\big)/(d(d+1))$. For $t=3$ this is $1/d^2$, equal to
$6/(d(d+1)(d+2))$ only for $d=2$. So the set is a 3-design exactly when $d=2$, where it is the octahedron of six
Pauli eigenstates; its order-4 frame potential $5/24$ exceeds $1/D_4=1/5$, so it is not a 4-design.
In each case the stated threshold is the design order.
\end{proof}

\paragraph{Exact gaps.} Table~\ref{tab:cert} lists the first strict gaps. Each value is certified in exact
Gaussian-rational arithmetic by a scalar dual $Y=\lambda I$ (positive semidefiniteness of $I\otimes Y-X$ verified by an
exact $LDL^\dagger$ factorization) and a matching primal Choi operator (an exact projector onto the corresponding
eigenspace, rescaled, with $\Tr_{\rm out}J=I$ verified exactly), so primal and dual values coincide.

\begin{table}[ht]
\centering\small
\caption{Optimal average global fidelity versus the universal value. ``Exact'': rational primal and dual
certificates; [Num]: floating-point semidefinite programming, not certified.}\label{tab:cert}
\begin{tabular}{@{}lccccl@{}}\toprule
Ensemble & design order & task & $P_\mu$ & $D_T/D_{T+k}$ & status\\\midrule
octahedron ($d=2$) & 3 & $1\to2,\,1\to3$ & $2/3,\ 1/2$ & $2/3,\ 1/2$ & exact\\
octahedron ($d=2$) & 3 & $1\to4$ & $5/12$ & $2/5$ & exact\\
complete MUBs ($d=4$) & 2 & $1\to2$ & $2/5$ & $2/5$ & exact\\
complete MUBs ($d=4$) & 2 & $1\to3$ & $1/4$ & $1/5$ & exact\\
2-qubit stabilizer states (60) & 3 & $1\to4$ & $1/8$ & $4/35$ & exact\\
complete MUBs ($d=3$) & 2 & $1\to3$ & $0.333333$ & $3/10$ & [Num]\\
qubit SIC (tetrahedron) & 2 & $1\to3$ & $0.555556$ & $1/2$ & [Num]\\\bottomrule
\end{tabular}
\end{table}

\section{Structured families: a contrast}\label{sec:contrast}
Theorem~\ref{thm:main} compares an ensemble with the \emph{universal} benchmark. A different question compares a
finite ensemble with the continuous family it is drawn from, for example a finite set of phase states with the full
torus of phase states, whose optimal cloning value $\Ptor$ was determined in~\cite[Thm.~II.3, Eq.~(16)]{Herasymenko2026}. There the
relation between moments and cloning is different, as shown in~\cite{Weldon2026}:

\begin{center}\small
\begin{tabular}{@{}p{0.42\linewidth}p{0.33\linewidth}p{0.17\linewidth}@{}}\toprule
Comparison & Matching the $(T+k)$-th moment & Reference\\\midrule
ensemble vs.\ universal value ($U(d)$) & necessary and sufficient & Theorem~\ref{thm:main}\\
finite phase codebook vs.\ torus value $\Ptor$ & sufficient in the proved cases; not necessary & \cite[Thm.~9, \S5.4]{Weldon2026}\\
finite orbit vs.\ continuous orbit of a subgroup $H\subsetneq U(d)$ & not sufficient in general & \cite[Thm.~31]{Weldon2026}\\\bottomrule
\end{tabular}
\end{center}

For phase codebooks, a key set $K\subset\mathbb Z_q$ that is $B_{T+k}$ (all $(T+k)$-fold sums distinct) reaches
$\Ptor$ in the proved cases, yet at $T=1$ there are key sets reaching $\Ptor$ without being $B_{1+k}$, for example
$\{0,1,13\}\subset\mathbb Z_{39}$ at $1\to3$~\cite[\S5.4]{Weldon2026}. For orbits of subgroups such as $SU(3)$ on
$\Sym^2\C^3$ and $SU(2)$ on spin~$1$, finite orbits matching the continuous orbit's moment can still be strictly
easier to clone~\cite[Thm.~31]{Weldon2026}. The proof of Theorem~\ref{thm:main} uses the full unitary group in an
essential way: it tests the identity~\eqref{eq:onesided} on \emph{every} coherent tensor power $\phi^{\otimes n}$, which
is possible because the universal cloner is built from the full symmetric projector.

\section{Discussion}\label{sec:discussion}
\paragraph{Relation to prior work.} The universal value and its uniqueness are due to Werner~\cite[Sec.~4]{Werner1998}.
Sufficiency of designs appears in different forms in~\cite{MVW2012,Kay2014}; in particular, Molina, Vidick and Watrous
observe that the four-state qubit SIC already attains $2/3$ for $1\to2$ cloning~\cite{MVW2012}, consistent with
Theorem~\ref{thm:main} since it is a 2-design. Numerical and semidefinite-programming approaches to cloning are
standard~\cite{Koniorczyk2013}. In measurement-based settings, Scott characterizes the optimal measurement-based
cloners~\cite[Thm.~20]{Scott2006} and Kim characterizes the ensembles of minimal quantumness~\cite{Kim2006};
Theorem~\ref{thm:main} is a coherent analogue, with a different objective and an unrestricted channel class. Searches of the literature (preprint servers
and two independent literature-search tools) did not locate a statement or proof of the necessity direction for
cloning; we would be grateful to learn of one.

\paragraph{What remains open.} Theorem~\ref{thm:main} decides when the universal value is optimal, not the optimum
itself. For stabilizer states it shows that, for average global fidelity, the universal cloner stops being optimal exactly
at output order four. The optimal channel beyond that point, the trace-distance formulation of~\cite{Herasymenko2026},
and the relation of the strict fidelity gap to the sample-complexity results of~\cite{Bansal2026} remain open. Two further questions are whether a quantitative version holds (a gap $P_\mu-D_T/D_{T+k}$
bounded below by a distance of $M_{T+k}$ from $\Pi/D_{T+k}$), and what the analogue is for worst-case fidelity.

\section*{Use of AI tools}
This work used several AI systems; the author directed the work and takes full responsibility for its content.
\textbf{Grok Bot} (xAI's agent platform, launched on the Grok~4.6 model) proposed the theorem and its proof and
produced the exact certificates. \textbf{ChatGPT} (OpenAI; GPT-6 Luna, the ChatGPT Free model) independently checked the proof line by line and
reviewed the draft.
\textbf{Perplexity} (Perplexity Pro with the GPT-6.1 Sol model) and \textbf{Consensus} carried out independent
prior-art searches; Perplexity also audited the citations. \textbf{Claude} (Anthropic; Claude Opus 5.5) audited the proof, reproduced the exact certificates
with independently generated state sets, verified every reference against arXiv or Crossref, and drafted the
manuscript.

\section*{Data and code availability}
The certificate scripts and their outputs (\texttt{paper2\_certificates.zip}, with a README and SHA-256 checksums) are
deposited with this note in the same Zenodo record, and in the folder \texttt{paper2/} of the repository
\url{https://github.com/Weldog/sidon-phase-cloning}. Verification requires Python $\ge3.10$ and \texttt{numpy}.

\section*{Author contributions and competing interests}
M.~Weldon is the sole author. He conceived and directed the project, read and approved the manuscript, and takes
responsibility for its content, with AI assistance as described above. The author declares no competing interests.
This work received no external funding.

\bibliographystyle{unsrtnat}
\bibliography{refs2}
\end{document}
